\section{Giới thiệu}

\subsection{Bối cảnh}

Mỗi năm, Trường Đại học Bách khoa – ĐHQG-HCM tiếp nhận hàng nghìn sinh viên mới, và mỗi sinh
viên trong suốt khoá học đều phải tra cứu quy định học vụ ít nhất vài chục lần: đăng ký tối
thiểu bao nhiêu tín chỉ, bị cảnh báo học tập trong trường hợp nào, điểm chữ quy đổi ra thang 4
thế nào, điều kiện xét tốt nghiệp gồm những gì, rút học phần trước thời điểm nào thì không bị
tính học phí.

Nguồn tra cứu chính thức là văn bản \emph{“Quy định về học vụ và đào tạo bậc đại học – Phiên
bản hợp nhất”} (số 285/ĐHBK-ĐT, ban hành ngày 28/7/2022). Đây là một văn bản hợp nhất từ bốn
nguồn: Thông tư 08/2021/TT-BGDĐT, Thông tư 17/2021/TT-BGDĐT, Quy định chung về học vụ và đào
tạo (QĐ 2931/QĐ-ĐHBK), và Quy định học vụ bậc đại học (QĐ 2933/QĐ-ĐHBK). Văn bản dài 57 trang,
gồm 10 chương, 42 điều và 15 phụ lục; mỗi đoạn đều kèm chú thích ở lề phải ghi rõ nội dung đó
bắt nguồn từ văn bản gốc nào.

Trong thực tế, rất ít sinh viên đọc hết 57 trang này. Họ hỏi bạn bè, hỏi trên các nhóm mạng xã
hội, hoặc hỏi trực tiếp cố vấn học tập và Phòng Đào tạo. Hệ quả là thông tin truyền miệng sai
lệch lan rộng (“học lại là bị khoá điểm cũ”, “nợ 3 môn là bị đuổi”), còn cán bộ phải trả lời đi
trả lời lại cùng một nhóm câu hỏi.

\subsection{Phát biểu bài toán}

Một trợ lý hỏi đáp tự động là giải pháp tự nhiên, nhưng dùng thẳng một mô hình ngôn ngữ lớn
(LLM) cho bài toán này là không chấp nhận được: mô hình sẽ trả lời trôi chảy kể cả khi nó không
biết, và với văn bản quy phạm thì một con số bịa (“tối đa 14 học kỳ” thay vì 12) nguy hiểm hơn
nhiều so với việc không trả lời. Bài toán vì vậy có bốn ràng buộc:

\begin{enumerate}
  \item \textbf{Chỉ trả lời dựa trên văn bản.} Mọi câu trả lời phải bắt nguồn từ nội dung thực
  sự có trong quy định, không suy đoán, không dùng kiến thức chung.
  \item \textbf{Truy vết được.} Người đọc phải kiểm chứng được ngay: mỗi ý kèm trích dẫn
  Điều/Khoản, và mở được đúng đoạn văn bản gốc cùng số trang trong PDF.
  \item \textbf{Biết nói “không biết”.} Khi văn bản không đề cập, hệ thống phải từ chối thay vì
  bịa.
  \item \textbf{Chạy được tại chỗ.} Dữ liệu học vụ là dữ liệu nội bộ; hệ thống phải chạy hoàn
  toàn trên hạ tầng của trường, không bắt buộc gửi câu hỏi ra dịch vụ bên ngoài.
\end{enumerate}

\subsection{Mục tiêu và phạm vi}

\textbf{Mục tiêu.} Xây dựng một hệ thống Retrieval-Augmented Generation (RAG) hoàn chỉnh, chạy
local, trả lời câu hỏi về quy định học vụ bằng tiếng Việt kèm trích dẫn kiểm chứng được; đồng
thời đo được chất lượng truy xuất một cách định lượng thay vì cảm tính.

\textbf{Phạm vi.} Đồ án tập trung vào một văn bản quy định học vụ hợp nhất, nhưng toàn bộ
pipeline được thiết kế đa tài liệu: thêm văn bản mới chỉ cần bỏ file vào thư mục nguồn và chạy
lại lệnh nạp dữ liệu. Phạm vi \emph{không} bao gồm xác thực người dùng, tích hợp với hệ thống
MyBK, và tra cứu dữ liệu cá nhân của sinh viên (điểm số, thời khoá biểu).

\subsection{Đóng góp chính}

So với một pipeline RAG mẫu, hệ thống có năm điểm được thiết kế riêng cho bài toán văn bản quy
phạm tiếng Việt:

\begin{itemize}
  \item \textbf{Phân mảnh bám cấu trúc pháp lý.} Đơn vị phân mảnh là Điều/Phụ lục chứ không
  phải cửa sổ trượt N ký tự; bảng biểu không bao giờ bị cắt đôi.
  \item \textbf{BM25 viết riêng cho tiếng Việt.} Chỉ mục theo âm tiết và bigram âm tiết, thay
  cho các mô hình BM25 sẵn có vốn stemming theo tiếng Anh.
  \item \textbf{Mở rộng thuật ngữ sinh viên → văn bản.} Bảng tra ánh xạ cách nói đời thường
  (“GPA”, “rớt môn”, “đuổi học”) và cả câu hỏi gõ không dấu sang đúng thuật ngữ trong quy định.
  \item \textbf{Rerank không cần cross-encoder.} Chấm lại ứng viên bằng các tín hiệu mà tìm
  kiếm vector bỏ qua, giữ nguyên yêu cầu không thêm mô hình nặng.
  \item \textbf{Soát trích dẫn tự động.} Mọi trích dẫn trong câu trả lời được đối chiếu với các
  Điều thực sự nằm trong ngữ cảnh; trích dẫn không có căn cứ bị đánh dấu cảnh báo.
\end{itemize}

Ngoài phần hệ thống, đồ án còn xây dựng bộ công cụ đánh giá gồm 30 câu hỏi vàng và một harness
đo \code{recall@k}, \code{hit@final}, \code{MRR}, cho phép bật/tắt từng kỹ thuật để đo đóng góp
riêng của nó.

\chapterbreak
